\chapter{Extended Calculations}
\label{appendix:calculations}

In this appendix, we provide detailed calculations that are relevant to the main text but were omitted for brevity.
The appendix is organized into separate sections that correspond to specific intermediate results used throughout the dissertation.

\section*{Total Transmit Power Constraint}
\label{calculation:transmit_power_constraint}

The following calculations derive the total transmit power as a function of the precoding matrix for both the \ac{SU}-\ac{MIMO} and \ac{MU}-\ac{MIMO} system models.
In both cases, the power allocation is assumed to be incorporated into the precoding matrix, and the data symbols are assumed to be zero-mean, mutually uncorrelated, and normalized, i.e.\ $\mathbb{E}\!\left[\mathbf{a}\mathbf{a}^H\right] = \mathbf{I}_{N_S}$.

Consider the \ac{SU}-\ac{MIMO} system model as described in Section~\ref{subsection:su_mimo_system_model}.
The total transmit power is constraint is given by:
\begin{equation*}
    \mathbb{E}\!\left[\|\mathbf{x}\|^2\right] = \mathrm{tr}\!\left(\mathbf{F}^H \mathbf{F}\right) \leq P_t.
\end{equation*}

\textbf{Calculation:}
\begin{fleqn}[2.5em]
\begin{subequations}
\begin{align*}
    P_t
    &\geq \mathbb{E}\!\left[ \left\|\mathbf{x}\right\|^2 \right]
    = \mathbb{E}\!\left[ \mathbf{x}^H \mathbf{x} \right]
    = \mathbb{E}\!\left[ \left(\mathbf{F} \mathbf{a}\right)^H \left(\mathbf{F} \mathbf{a}\right) \right]
    = \mathbb{E}\!\left[ \mathbf{a}^H \mathbf{F}^H \mathbf{F} \mathbf{a} \right] \\
    &= \mathbb{E}\!\left[ \mathrm{tr}\!\left( \mathbf{a}^H \mathbf{F}^H \mathbf{F} \mathbf{a} \right) \right]
    = \mathbb{E}\!\left[ \mathrm{tr}\!\left( \mathbf{F}^H \mathbf{F} \mathbf{a} \mathbf{a}^H \right) \right] \\
    &= \mathrm{tr}\!\left( \mathbb{E}\!\left[ \mathbf{F}^H \mathbf{F} \mathbf{a} \mathbf{a}^H \right] \right)
    = \mathrm{tr}\!\left( \mathbf{F}^H \mathbf{F} \, \mathbb{E}\!\left[ \mathbf{a} \mathbf{a}^H \right] \right) \\
    &= \mathrm{tr}\!\left( \mathbf{F}^H \mathbf{F} \, \mathbf{I}_{N_S} \right)
    = \mathrm{tr}\!\left( \mathbf{F}^H \mathbf{F} \right)
\end{align*}
\end{subequations}
\end{fleqn}

Consider the \ac{MU}-\ac{MIMO} system model as described in Section~\ref{subsection:mu_mimo_system_model}.
The total transmit power is constraint is given by:
\begin{equation*}
    \mathbb{E}\!\left[\|\mathbf{x}\|^2\right] = \sum_{k=1}^{K} \mathrm{tr}\!\left(\mathbf{F}_k^H \mathbf{F}_k\right) \leq P_t.
\end{equation*}

\textbf{Calculation:}
\begin{fleqn}[2.5em]
\begin{subequations}
\begin{align*}
    P_t
    &\geq \mathbb{E}\!\left[ \left\|\mathbf{x}\right\|^2 \right]
    = \mathbb{E}\!\left[ \mathbf{x}^H \mathbf{x} \right]
    = \mathbb{E}\!\left[ \left(\sum_{k=1}^{K} \mathbf{F}_k \mathbf{a}_k\right)^H \left(\sum_{k'=1}^{K} \mathbf{F}_{k'} \mathbf{a}_{k'}\right) \right] \\
    &= \mathbb{E}\!\left[ \sum_{k=1}^{K} \sum_{k'=1}^{K} \left(\mathbf{F}_k \mathbf{a}_k \right)^H \mathbf{F}_{k'} \mathbf{a}_{k'} \right]
    = \mathbb{E}\!\left[ \sum_{k=1}^{K} \sum_{k'=1}^{K} \mathbf{a}_k^H \mathbf{F}_k^H \mathbf{F}_{k'} \mathbf{a}_{k'} \right] \\
    &= \mathbb{E}\!\left[ \sum_{k=1}^{K} \sum_{k'=1}^{K} \text{Tr}\!\left( \mathbf{a}_k^H \mathbf{F}_k^H \mathbf{F}_{k'} \mathbf{a}_{k'} \right) \right]
    = \mathbb{E}\!\left[ \sum_{k=1}^{K} \sum_{k'=1}^{K} \text{Tr}\!\left( \mathbf{F}_k^H \mathbf{F}_{k'} \, \mathbf{a}_{k'} \mathbf{a}_k^H \right) \right] \\
    &= \sum_{k=1}^{K} \sum_{k'=1}^{K} \text{Tr}\!\left( \mathbf{F}_k^H \mathbf{F}_{k'} \, \mathbb{E}\!\left[ \mathbf{a}_{k'} \mathbf{a}_k^H \right] \right) \\
    &= \sum_{k=1}^{K} \sum_{k'=1}^{K} \text{Tr}\!\left( \mathbf{F}_k^H \mathbf{F}_{k'} \, \mathbf{I}_{N_S} \, \delta_{kk'} \right) \\
    &= \sum_{k=1}^{K} \text{Tr}\!\left( \mathbf{F}_k^H \mathbf{F}_{k} \right)
\end{align*}
\end{subequations}
\end{fleqn}

\section*{Equivalent SU-MIMO Channel Representation}
\label{calculation:su_mimo_equivalent_channel}

Applying the singular value decomposition of the channel matrix together with a change of variables yields an equivalent representation of the \ac{SU}-\ac{MIMO} channel. The following calculation shows how the original channel input-output relation is transformed into a set of parallel scalar Gaussian subchannels.

\textbf{Calculation:}
\begin{fleqn}[2.5em]
\begin{subequations}
\begin{align*}
    & \mathbf{y} = \mathbf{H} \mathbf{x} + \mathbf{n} \\
    & \Leftrightarrow \quad \mathbf{y} = \mathbf{U} \mathbf{\Sigma} \mathbf{V}^H \, \mathbf{x} + \mathbf{n} \\
    & \Leftrightarrow \quad \mathbf{U} \tilde{\mathbf{y}} = \mathbf{U} \mathbf{\Sigma} \mathbf{V}^H \, \mathbf{V} \tilde{\mathbf{x}} + \mathbf{n} \\
    & \Leftrightarrow \quad \tilde{\mathbf{y}} = \mathbf{U}^H \mathbf{U} \mathbf{\Sigma} \mathbf{V}^H \mathbf{V} \, \tilde{\mathbf{x}} + \mathbf{U}^H \mathbf{n} \\
    & \Leftrightarrow \quad \tilde{\mathbf{y}} = \mathbf{\Sigma} \tilde{\mathbf{x}} + \tilde{\mathbf{n}}.
\end{align*}
\end{subequations}
\end{fleqn}

\section*{Mutual Information under Optimal Precoder and Combiner}
\label{calculation:su_mimo_optimal_precoder_mi}

The following calculation verifies that the \ac{MI} between the transmitted data symbols $\mathbf{a}$ and the decision variables $\mathbf{z}$ equals the \ac{MI} between the channel input $\mathbf{x}$ and the channel output $\mathbf{y}$ when the optimal precoder and combiner from~\eqref{eq:su_mimo_optimal_precoder_and_combiner} are employed.
The result is obtained by substituting $\mathbf{F} = \mathbf{V}_{N_s}\sqrt{\mathbf{P}}$ and $\mathbf{G} = \mathbf{U}_{N_s}^H$ into~\eqref{eq:mi_az_b} and replacing the channel matrix by its \ac{SVD}~\eqref{eq:svd_channel_matrix}.

\textbf{Calculation:}
\begin{fleqn}[1.5em]
\begin{subequations}
\begin{align*}
    &I(\mathbf{a}; \mathbf{z}) \\
    &= \log_2\!\left(\det\!\left(\mathbf{I}_{N_s} + \tfrac{1}{N_0} (\mathbf{U}_{N_s}^H \mathbf{U}_{N_s})^{-1} \, (\mathbf{U}_{N_s}^H (\mathbf{U} \mathbf{\Sigma} \mathbf{V}^H) \mathbf{V}_{N_s} \sqrt{\mathbf{P}}) \, (\mathbf{U}_{N_s}^H (\mathbf{U} \mathbf{\Sigma} \mathbf{V}^H) \mathbf{V}_{N_s} \sqrt{\mathbf{P}})^H\right)\right) \\
    &= \log_2\!\left(\det\!\left(\mathbf{I}_{N_s} + \tfrac{1}{N_0} \mathbf{I}_{N_s} \, (\mathbf{U}_{N_s}^H \mathbf{U} \mathbf{\Sigma} \mathbf{V}^H \mathbf{V}_{N_s} \sqrt{\mathbf{P}}) \, (\mathbf{U}_{N_s}^H \mathbf{U} \mathbf{\Sigma} \mathbf{V}^H \mathbf{V}_{N_s} \sqrt{\mathbf{P}})^H\right)\right) \\
    &= \log_2\!\left(\det\!\left(\mathbf{I}_{N_s} + \tfrac{1}{N_0} \, \left(\begin{bmatrix} \mathbf{I}_{N_s} & \mathbf{0} \end{bmatrix} \, \begin{bmatrix} \mathbf{\Sigma}_{R_H} & \mathbf{0} \\ \mathbf{0} & \mathbf{0} \end{bmatrix} \, \begin{bmatrix} \mathbf{I}_{N_s} \\ \mathbf{0} \end{bmatrix} \sqrt{\mathbf{P}}\right) \, \left(\begin{bmatrix} \mathbf{I}_{N_s} & \mathbf{0} \end{bmatrix} \, \begin{bmatrix} \mathbf{\Sigma}_{R_H} & \mathbf{0} \\ \mathbf{0} & \mathbf{0} \end{bmatrix} \, \begin{bmatrix} \mathbf{I}_{N_s} \\ \mathbf{0} \end{bmatrix} \sqrt{\mathbf{P}}\right)^H\right)\right)  \\
    &= \log_2\!\left(\det\!\left(\mathbf{I}_{N_s} + \tfrac{1}{N_0} \, (\mathbf{\Sigma}_{N_s} \sqrt{\mathbf{P}}) \, (\mathbf{\Sigma}_{N_s} \sqrt{\mathbf{P}})^H\right)\right) \\
    &= \log_2\!\left(\det\!\left(\mathbf{I}_{N_s} + \tfrac{1}{N_0} \, \mathbf{\Sigma}_{N_s}^2 \, \mathbf{P}\right)\right) \\
    &= \log_2\!\left( \prod_{i=1}^{N_s} \left(1 + \frac{\sigma^2_i}{N_0}\, p_i \right) \right)  \\
    &= \sum_{i=1}^{N_s} \log_2\!\left(1 + \frac{\sigma^2_i}{N_0}\, p_i \right) \\
    %&= I(\mathbf{x}; \mathbf{x}) \qquad \text{for} \quad p_i = \mathbb{E}\left[\left\|x\right\|^2\right]
\end{align*}
\end{subequations}
\end{fleqn}
